%!TEX root = ../../main.tex
교전은 킬을 씨앗으로 시간 경계 $G$와 공간 경계 $D$로 묶은 뒤 존재 게이트 $R$, $B$, $M$으로 걸러 얻는다. 이 장은 $G$와 $D$를 코퍼스에서 추정하고, $R$과 $B$를 게임 자체가 적용하는 규칙에 고정하며, 나머지 상수를 각각의 근거 또는 ``측정되지 않았음''과 함께 표로 남기는 규칙(검출기 v3.3)을 정의한다. 이어서 규모 계급과 코호트, 결과 시점 규칙, 킬에 조건화된 정의의 분석 대상, 정의 상수의 민감도를 다룬다. 정의 상수의 추정치와 구간은 상수 추정 기록(부록 \ref{app:repro})에서 옮긴 값이고, Baek과 Kwon \cite{baek2026killconditioned}의 상수(킬 간격 \cogGap{}초, 군집 지름 \cogDiam{} 단위, 게이트 \cogGateR{} 단위·\cogLead{}초)는 제 \ref{sec:eng-temporal} 절과 \ref{sec:eng-gate} 절에서 비교 기준으로 쓴다. 교전 승자 라벨(\texttt{market\_event}) 아래에서 측정된 예측 민감도는 그 라벨의 결과로 인용한다.
